\chapter{Introduction}\label{ch:intro}

\ac{} computes anharmonic infrared spectra of molecules, in absorption and
in emission. It covers the fundamentals and the overtone, combination and
difference bands, and the hot bands of vibrationally excited molecules. It
is designed for large molecules, such as polycyclic aromatic hydrocarbons
(PAHs), and for spectra of molecules with a given internal energy or
temperature. The input is the same quartic force field and second-order
dipole that an anharmonic frequency calculation with an electronic
structure code provides; Gaussian output can be converted directly
(Chapter~\ref{ch:tut-gaussian}).

\section{What it does, compared with VPT2}

A VPT2 calculation, such as \texttt{freq=anharmonic} in Gaussian, gives
anharmonic fundamentals, overtones and combinations from the ground state.
Resonances are treated by deperturbation and small variational blocks.
\ac{} proceeds in three steps.
\begin{enumerate}
\item The \emph{same perturbation theory}, in its canonical (Van Vleck)
  form, removes the anharmonic terms except those that couple nearly
  degenerate states: the Fermi and Darling--Dennison resonances. This gives
  the deperturbed anharmonic constants, the effective resonant couplings
  and the dipole transformed to second order \notes{§3--6}.
\item For each starting state, the states connected by the resonant
  couplings are collected into \emph{polyads}. These are finite blocks of
  harmonic states, truncated by an energy and coupling criterion
  \notes{§11.3}.
\item The polyad Hamiltonians are \emph{diagonalised}, or treated
  spectrally for very large polyads. The transitions between the resulting
  states are accumulated into histograms, binned in wavenumber and in the
  energy of the starting state.
\end{enumerate}
Since any starting state can be used, not only the ground state, averaging
over a thermal or microcanonical sample of starting states gives band
positions, widths and shapes as a function of temperature or internal
energy. A plain VPT2 calculation does not give this.

\section{When to use it}

\begin{itemize}
\item Spectra of vibrationally hot molecules, for example the infrared
  emission of PAHs heated by the absorption of a UV photon.
\item Large molecules with many accidental resonances, where the
  variational treatment of the polyads matters.
\item Temperature-dependent band profiles.
\end{itemize}
For the fundamentals of a small molecule at 0~K, a VPT2 calculation is
usually enough; \ac{} then reproduces it, as a check
(Chapter~\ref{ch:tut-gaussian}).

\section{How to read this manual}

Chapters~\ref{ch:install} and~\ref{ch:workflow} explain how to build the
code and how a calculation is organised. Chapters~\ref{ch:tut-pyrene}
to~\ref{ch:tut-review} are tutorials with complete, reproducible
examples. Chapters~\ref{ch:inputs} and~\ref{ch:outputs} are the
reference for the input and output files. Chapter~\ref{ch:trouble}
explains the warnings and errors.

The theory behind the code, with all the derivations, is in the companion
notes \emph{Anharmonic oscillations with the Van Vleck method}, cited here
as ``theory notes''. The method is described in~\cite{mulas2018}; the public release is on
SourceForge~\cite{anharmonicasf}.
